\documentclass[11pt]{article}

\usepackage[a4paper,margin=1in]{geometry}
\usepackage{amsmath,amssymb}
\usepackage{hyperref}

\hypersetup{
    colorlinks=true,
    linkcolor=blue,
    citecolor=blue,
    urlcolor=blue
}

\title{Why Transport Maps Cannot Split Mass}
\author{}
\date{}

\begin{document}

\maketitle

A transport map cannot split mass because it is a deterministic function.
If
\[
T:X\to Y
\]
is a map, then every source point $x$ is sent to exactly one target point
$T(x)$. Thus all mass located at $x$ must move together to the single point
$T(x)$.

In measure notation, the transported measure is
\[
T_\#\mu(A)=\mu(T^{-1}(A)).
\]
This formula has no mechanism for sending part of the mass at $x$ to one point
and another part to a different point.

\section*{Example}

Let
\[
\mu=\delta_0,
\qquad
\nu=\frac12\delta_{-1}+\frac12\delta_1.
\]
No transport map $T$ can satisfy $T_\#\mu=\nu$, because
\[
T_\#\delta_0=\delta_{T(0)}.
\]
This is always a single Dirac mass, not a mixture of two Dirac masses.

A transport plan can split the mass:
\[
\pi
=
\frac12\delta_{(0,-1)}
+
\frac12\delta_{(0,1)}.
\]
This plan sends half the mass at $0$ to $-1$ and half to $1$.

\end{document}
